% ============================================================
%  C: Como Programar -- 6ª Edição | Deitel & Deitel
%  Capítulo 7 -- Ponteiros em C
%  EXERCÍCIOS (7.7 -- 7.31)
%  Abordagem incremental: Caps. 1 + 2 + 3 + 4 + 5 + 6 + 7
% ============================================================
\documentclass[12pt,a4paper]{article}
\usepackage[utf8]{inputenc}
\usepackage[T1]{fontenc}
\usepackage[brazil]{babel}
\usepackage{geometry}
\geometry{top=2.5cm,bottom=2.5cm,left=3cm,right=3cm}
\usepackage{parskip}\setlength{\parskip}{6pt}\setlength{\parindent}{0pt}
\usepackage{xcolor,tcolorbox,enumitem,listings,fancyhdr,titlesec,hyperref,booktabs,array,amsmath}
\tcbuselibrary{skins,breakable}

\definecolor{deitelblue}{RGB}{0,70,127}
\definecolor{deitelgray}{RGB}{240,242,245}
\definecolor{deitelgreen}{RGB}{0,110,60}
\definecolor{deitellightblue}{RGB}{220,235,250}
\definecolor{deitelred}{RGB}{160,20,20}
\definecolor{deitellorange}{RGB}{200,90,0}

\newtcolorbox{boaprática}[1][]{colback=deitellightblue,colframe=deitelblue,
  fonttitle=\bfseries\small,title={Boa prática de programação},breakable,#1}
\newtcolorbox{erroco}[1][]{colback=red!5,colframe=deitelred,
  fonttitle=\bfseries\small,title={Erro comum de programação},breakable,#1}
\newtcolorbox{respostabox}[1][]{colback=deitelgray,colframe=deitelblue!60,
  fonttitle=\bfseries\small\color{deitelblue},title={Resposta detalhada},breakable,#1}
\newtcolorbox{enunciadoBox}[1][]{colback=white,colframe=deitelblue,
  fonttitle=\bfseries,title={#1},breakable}
\newtcolorbox{saida}[1][]{colback=black!88,colframe=deitelblue,
  fontupper=\ttfamily\small\color{green!70},
  title={\small\color{white}Saída do programa},breakable,#1}

\setlist[enumerate,1]{label=\textbf{\alph*)},leftmargin=2em}
\setlist[itemize]{leftmargin=2em}

\lstset{
  basicstyle=\ttfamily\small,backgroundcolor=\color{deitelgray},
  frame=single,framerule=0.4pt,rulecolor=\color{deitelblue!40},
  breaklines=true,showstringspaces=false,
  keywordstyle=\color{deitelblue}\bfseries,
  commentstyle=\color{deitelgreen}\itshape,
  stringstyle=\color{deitelred},
  language=C,numbers=left,numberstyle=\tiny\color{gray},
  xleftmargin=18pt,xrightmargin=5pt,stepnumber=1,
}

\pagestyle{fancy}\fancyhf{}
\fancyhead[L]{\small\color{deitelblue}\textbf{C: Como Programar -- 6ª ed. | Deitel \& Deitel}}
\fancyhead[R]{\small\color{deitelblue}Cap. 7 -- Exercícios}
\fancyfoot[C]{\small\thepage}
\renewcommand{\headrulewidth}{0.4pt}

\titleformat{\section}[block]{\large\bfseries\color{deitelblue}}{\thesection.}{0.5em}{}[\titlerule]
\titleformat{\subsection}[block]{\normalsize\bfseries\color{deitelblue!80}}{}{}{}

\hypersetup{colorlinks=true,linkcolor=deitelblue}

\begin{document}

\begin{titlepage}
  \centering\vspace*{2cm}
  {\color{deitelblue}\rule{\linewidth}{2pt}}\\[0.4cm]
  {\huge\bfseries\color{deitelblue}C: Como Programar}\\[0.2cm]
  {\large\color{deitelblue}6ª Edição $\cdot$ Paul Deitel \& Harvey Deitel}\\[0.2cm]
  {\color{deitelblue}\rule{\linewidth}{2pt}}\\[1cm]
  {\LARGE\bfseries Capítulo 7}\\[0.3cm]
  {\Large\color{deitelblue!80}Ponteiros em C}\\[0.8cm]
  \begin{tcolorbox}[colback=deitellightblue,colframe=deitelblue,width=0.85\linewidth]
    \centering{\large\bfseries Exercícios (7.7--7.31)}\\[0.2cm]
    {\normalsize Resolvidos passo a passo, com código completo}
  \end{tcolorbox}
\end{titlepage}

\tableofcontents\newpage

% ============================================================
\section{Exercício 7.7 -- Preenchimento de Lacunas}
% ============================================================

\begin{respostabox}
\begin{enumerate}[label=\textbf{\alph*)}]
\item O operador \textbf{\&} (endereço, ou \textit{address-of}) retorna o local na memória onde o operando está armazenado.
\item O operador \textbf{*} (desreferência, ou \textit{indireção}) retorna o valor do objeto apontado pelo operando.
\item Para simular chamada por referência ao passar variável não-array, é necessário passar o(a) \textbf{endereço} da variável (usando \texttt{\&variavel}).
\end{enumerate}
\end{respostabox}

% ============================================================
\section{Exercício 7.8 -- Verdadeiro ou Falso}
% ============================================================

\begin{respostabox}
\begin{enumerate}[label=\textbf{\alph*)}]

\item \textcolor{deitelgreen}{\textbf{Verdadeiro.}} Comparar ponteiros que apontam
para arrays diferentes \textit{não tem significado prático}. As comparações só
fazem sentido para ponteiros dentro do mesmo array, onde pode-se verificar
qual elemento vem antes/depois.

\item \textcolor{deitelred}{\textbf{Falso.}} Embora o nome do array atue como
ponteiro para o primeiro elemento, ele é um \textbf{ponteiro constante} -- não pode
ser modificado. Tentar \texttt{++arr} ou \texttt{arr = outroPonteiro} gera erro de
compilação. Já uma variável-ponteiro pode ser livremente modificada.
\end{enumerate}
\end{respostabox}

% ============================================================
\section{Exercício 7.9 -- Array \texttt{values}}
% ============================================================

\begin{respostabox}
\textbf{Considere:} \texttt{unsigned int} = 2 bytes; endereço inicial = 1002500;
\texttt{\#define SIZE 5}.

\begin{enumerate}[label=\textbf{\alph*)}]
\item \begin{lstlisting}
unsigned int values[ SIZE ] = { 2, 4, 6, 8, 10 };
\end{lstlisting}

\item \begin{lstlisting}
unsigned int *vPtr;
\end{lstlisting}

\item \begin{lstlisting}
for ( i = 0; i < SIZE; ++i )
    printf( "%u ", values[ i ] );
\end{lstlisting}

\item \begin{lstlisting}
vPtr = values;
vPtr = &values[ 0 ];
\end{lstlisting}

\item \begin{lstlisting}
for ( i = 0; i < SIZE; ++i )
    printf( "%u ", *( vPtr + i ) );
\end{lstlisting}

\item \begin{lstlisting}
for ( i = 0; i < SIZE; ++i )
    printf( "%u ", *( values + i ) );
\end{lstlisting}

\item \begin{lstlisting}
for ( i = 0; i < SIZE; ++i )
    printf( "%u ", vPtr[ i ] );
\end{lstlisting}

\item \textbf{Elemento 5 (último, índice 4):}
\begin{lstlisting}
values[ 4 ]        /* notacao de subscrito de array     */
*( values + 4 )    /* nome do array como ponteiro       */
vPtr[ 4 ]          /* notacao de subscrito de ponteiro  */
*( vPtr + 4 )      /* ponteiro/deslocamento             */
\end{lstlisting}

\item \textbf{Endereço de \texttt{vPtr + 3}:} $1002500 + 3 \times 2 = \mathbf{1002506}$;
valor: \texttt{values[3] = 8}.

\item \textbf{Endereço após \texttt{vPtr -= 4} (\texttt{vPtr} aponta para \texttt{values[4]}):}
\begin{itemize}[noitemsep]
  \item Endereço de \texttt{values[4]}: $1002500 + 4 \times 2 = 1002508$
  \item Após \texttt{vPtr -= 4}: $1002508 - 4 \times 2 = \mathbf{1002500}$
\end{itemize}
Esse é o endereço de \texttt{values[0]}, valor: \textbf{2}.
\end{enumerate}
\end{respostabox}

\newpage

% ============================================================
\section{Exercício 7.10 -- \texttt{long} e ponteiros}
% ============================================================

\begin{respostabox}
\begin{enumerate}[label=\textbf{\alph*)}]
\item \texttt{long *lPtr;}
\item \texttt{lPtr = \&value1;}
\item \texttt{printf( "\%ld\textbackslash n", *lPtr );}
\item \texttt{value2 = *lPtr;}
\item \texttt{printf( "\%ld\textbackslash n", value2 );}
\item \texttt{printf( "\%p\textbackslash n", \&value1 );}
\item \texttt{printf( "\%p\textbackslash n", lPtr );} -- \textbf{Sim}, é o mesmo
endereço de \texttt{value1}, pois \texttt{lPtr} foi inicializado em (b).
\end{enumerate}
\end{respostabox}

% ============================================================
\section{Exercício 7.11 -- Cabeçalhos e Protótipos}
% ============================================================

\begin{respostabox}
\begin{enumerate}[label=\textbf{\alph*)}]
\item \textbf{Cabeçalho de \texttt{zero}:}
\begin{lstlisting}
void zero( long bigIntegers[] )
\end{lstlisting}

\item \textbf{Protótipo de \texttt{zero}:}
\begin{lstlisting}
void zero( long bigIntegers[] );
\end{lstlisting}

\item \textbf{Cabeçalho de \texttt{add1AndSum}:}
\begin{lstlisting}
int add1AndSum( int oneTooSmall[] )
\end{lstlisting}

\item \textbf{Protótipo:}
\begin{lstlisting}
int add1AndSum( int oneTooSmall[] );
\end{lstlisting}
\end{enumerate}

\textit{Nota:} em C, \texttt{tipo nome[]} em parâmetro de função é equivalente a
\texttt{tipo *nome}. Ambos significam ``ponteiro para o tipo''.
\end{respostabox}

% ============================================================
\section{Exercício 7.12 -- Embaralhamento + Pôquer (5 cartas)}
% ============================================================

\begin{respostabox}
\textbf{Estratégia:} estender o programa de embaralhar/distribuir cartas para
distribuir mãos de 5 cartas e detectar combinações de pôquer.

\begin{lstlisting}
/* Ex7.12: poker.c -- esqueleto com funcoes principais */
#include <stdio.h>
#include <stdlib.h>
#include <time.h>
#define CARTAS 52
#define MAO 5

/* Tipos de combinacoes */
int temPar          ( const int mao[] );
int temDoisPares    ( const int mao[] );
int temTrinca       ( const int mao[] );
int temQuadra       ( const int mao[] );
int temFlush        ( const int mao[] );
int temStraight     ( const int mao[] );

/* As cartas vem representadas por inteiros 0..51:
   valor = carta % 13   (0=A, 1=2, ..., 12=K)
   naipe = carta / 13   (0=Ouros, 1=Espadas, 2=Copas, 3=Paus)  */

/* Conta frequencias de cada valor (0..12) na mao */
void contaValores( const int mao[], int freq[] )
{
    int i;
    for ( i = 0; i < 13; ++i ) freq[ i ] = 0;
    for ( i = 0; i < MAO; ++i ) {
        ++freq[ mao[ i ] % 13 ];
    }
}

int temPar( const int mao[] )
{
    int freq[ 13 ], i;
    contaValores( mao, freq );
    for ( i = 0; i < 13; ++i ) {
        if ( freq[ i ] == 2 ) return 1;
    }
    return 0;
}

int temDoisPares( const int mao[] )
{
    int freq[ 13 ], i, pares = 0;
    contaValores( mao, freq );
    for ( i = 0; i < 13; ++i ) {
        if ( freq[ i ] == 2 ) ++pares;
    }
    return pares == 2;
}

int temTrinca( const int mao[] )
{
    int freq[ 13 ], i;
    contaValores( mao, freq );
    for ( i = 0; i < 13; ++i ) {
        if ( freq[ i ] == 3 ) return 1;
    }
    return 0;
}

int temQuadra( const int mao[] )
{
    int freq[ 13 ], i;
    contaValores( mao, freq );
    for ( i = 0; i < 13; ++i ) {
        if ( freq[ i ] == 4 ) return 1;
    }
    return 0;
}

int temFlush( const int mao[] )
{
    int naipe = mao[ 0 ] / 13;
    int i;
    for ( i = 1; i < MAO; ++i ) {
        if ( mao[ i ] / 13 != naipe ) return 0;
    }
    return 1;
}

int temStraight( const int mao[] )
{
    int valores[ MAO ];
    int i, j, temp;

    /* Extrai e ordena os valores */
    for ( i = 0; i < MAO; ++i ) valores[ i ] = mao[ i ] % 13;
    for ( i = 0; i < MAO - 1; ++i ) {
        for ( j = 0; j < MAO - 1 - i; ++j ) {
            if ( valores[ j ] > valores[ j + 1 ] ) {
                temp = valores[ j ];
                valores[ j ] = valores[ j + 1 ];
                valores[ j + 1 ] = temp;
            }
        }
    }

    /* Verifica se sao consecutivos */
    for ( i = 0; i < MAO - 1; ++i ) {
        if ( valores[ i + 1 ] != valores[ i ] + 1 ) return 0;
    }
    return 1;
}
\end{lstlisting}
\end{respostabox}

\newpage

% ============================================================
\section{Exercício 7.13 -- Duas mãos de pôquer}
% ============================================================

\begin{respostabox}
\textbf{Estratégia:} distribuir 10 cartas (2 mãos de 5), avaliar cada mão e
determinar a melhor.

\begin{lstlisting}
/* Ex7.13: comparar_maos.c */
#include <stdio.h>
#include <stdlib.h>
#include <time.h>
#define CARTAS 52
#define MAO 5

/* Inclui funcoes de Ex7.12: temPar, temFlush, ... */

/* Retorna pontuacao da mao (maior = melhor) */
int avaliaMao( const int mao[] )
{
    if ( temFlush( mao ) && temStraight( mao ) ) return 8;  /* Straight Flush */
    if ( temQuadra( mao ) )                       return 7;  /* Quadra */
    if ( temTrinca( mao ) && temPar( mao ) )      return 6;  /* Full House */
    if ( temFlush( mao ) )                        return 5;  /* Flush */
    if ( temStraight( mao ) )                     return 4;  /* Sequencia */
    if ( temTrinca( mao ) )                       return 3;  /* Trinca */
    if ( temDoisPares( mao ) )                    return 2;  /* Dois pares */
    if ( temPar( mao ) )                          return 1;  /* Par */
    return 0;  /* Carta alta */
}

const char *nomeMao( int valor )
{
    switch ( valor ) {
        case 8: return "Straight Flush";
        case 7: return "Quadra";
        case 6: return "Full House";
        case 5: return "Flush";
        case 4: return "Sequencia";
        case 3: return "Trinca";
        case 2: return "Dois pares";
        case 1: return "Par";
        default: return "Carta alta";
    }
}

int main( void )
{
    int baralho[ CARTAS ];
    int mao1[ MAO ], mao2[ MAO ];
    int i, temp, idx;
    int v1, v2;

    srand( time( NULL ) );

    /* Inicializa e embaralha (Fisher-Yates) */
    for ( i = 0; i < CARTAS; ++i ) baralho[ i ] = i;
    for ( i = CARTAS - 1; i > 0; --i ) {
        idx = rand() % ( i + 1 );
        temp = baralho[ i ];
        baralho[ i ] = baralho[ idx ];
        baralho[ idx ] = temp;
    }

    /* Distribui 5 para cada mao */
    for ( i = 0; i < MAO; ++i ) {
        mao1[ i ] = baralho[ i ];
        mao2[ i ] = baralho[ i + MAO ];
    }

    v1 = avaliaMao( mao1 );
    v2 = avaliaMao( mao2 );

    printf( "Mao 1: %s\n", nomeMao( v1 ) );
    printf( "Mao 2: %s\n", nomeMao( v2 ) );

    if ( v1 > v2 ) printf( "Mao 1 vence!\n" );
    else if ( v2 > v1 ) printf( "Mao 2 vence!\n" );
    else printf( "Empate na categoria.\n" );

    return 0;
}
\end{lstlisting}
\end{respostabox}

\newpage

% ============================================================
\section{Exercícios 7.14 e 7.15 -- Carteador e Jogo Completo}
% ============================================================

\begin{respostabox}
\textbf{7.14 -- Simular o carteador trocando cartas:}

Após avaliar a mão do carteador, decida quantas cartas trocar:
\begin{itemize}[noitemsep]
  \item Se tem quadra ou melhor: \textit{não troca} (mão muito boa).
  \item Se tem trinca: troca 2 cartas restantes.
  \item Se tem dois pares: troca 1 carta (a do meio).
  \item Se tem um par: troca 3 cartas.
  \item Sem par algum: troca 3 cartas (mantém as 2 mais altas).
\end{itemize}

\textbf{7.15 -- Jogo completo (jogador vs.\ computador):}

\begin{lstlisting}
/* Esqueleto */
1. Embaralhe.
2. Distribua 5 cartas para o jogador e 5 para o carteador.
3. Mostre as cartas do jogador; carteador fica oculto.
4. Pergunte ao jogador quais cartas trocar (indices).
5. Substitua cartas do jogador.
6. O carteador escolhe automaticamente (regras do 7.14)
   quais trocar e troca.
7. Compare maos e exiba vencedor.
\end{lstlisting}

\textbf{Reflexão:} estes exercícios mostram como ponteiros e arrays compõem
sistemas complexos. O jogo de pôquer envolve combinatória, estratégia (do
carteador) e interatividade. Refinamentos sucessivos permitem ajustar a
heurística do computador.
\end{respostabox}

% ============================================================
\section{Exercício 7.16 -- Embaralhamento Eficiente}
% ============================================================

\begin{respostabox}
\textbf{Problema do algoritmo da Figura 7.24:} ele escolhe linha/coluna aleatórios
e \textit{tenta} colocar a carta -- se a posição já está ocupada, tenta outra,
podendo iterar indefinidamente (\textbf{adiamento indefinido}).

\textbf{Solução eficiente (\textit{Fisher-Yates Shuffle}):}

\begin{lstlisting}
/* Ex7.16: embaralhar eficiente -- O(n) */
void embaralhaEficiente( int baralho[][ 13 ] )
{
    int i, j, ni, nj, temp;
    /* Inicializa: posicao (linha, coluna) contem a carta original */

    /* Percorre todas as 52 posicoes */
    for ( i = 0; i < 4; ++i ) {
        for ( j = 0; j < 13; ++j ) {
            /* Escolhe posicao aleatoria QUALQUER do baralho */
            ni = rand() % 4;
            nj = rand() % 13;
            /* Troca */
            temp = baralho[ i ][ j ];
            baralho[ i ][ j ] = baralho[ ni ][ nj ];
            baralho[ ni ][ nj ] = temp;
        }
    }
}
\end{lstlisting}

\textbf{Vantagem:} cada elemento é tocado exatamente uma vez ($O(n)$). Não há
risco de adiamento indefinido. Esta é a essência do algoritmo de Fisher-Yates,
usado em todos os geradores profissionais de embaralhamento.

\textbf{Distribuição eficiente:} em vez de procurar a carta 1, depois a 2, etc.,
basta percorrer o array sequencialmente. Após embaralhar, as cartas \textit{já estão}
em ordem aleatória -- a posição no array \textit{é} a ordem de distribuição.
\end{respostabox}

\newpage

% ============================================================
\section{Exercício 7.17 -- A Tartaruga e a Lebre}
% ============================================================

\begin{respostabox}
\begin{lstlisting}
/* Ex7.17: tartaruga_lebre.c
   Simulacao classica da corrida com geracao aleatoria */
#include <stdio.h>
#include <stdlib.h>
#include <time.h>
#define PISTA 70

void moveTartaruga( int *posicao );
void moveLebre    ( int *posicao );
void exibeCorrida ( int tartaruga, int lebre );

int main( void )
{
    int tartaruga = 1, lebre = 1;
    int ticks = 0;

    srand( time( NULL ) );

    printf( "BANG !!!!!\nE LA VAO ELES !!!!!\n" );

    while ( tartaruga < PISTA && lebre < PISTA ) {
        moveTartaruga( &tartaruga );
        moveLebre( &lebre );
        exibeCorrida( tartaruga, lebre );
        ++ticks;
    }

    printf( "\n" );
    if ( tartaruga >= PISTA && lebre >= PISTA ) {
        printf( "E um empate.\n" );
    }
    else if ( tartaruga >= PISTA ) {
        printf( "TARTARUGA VENCE!!! E ISSO AI!!!\n" );
    }
    else {
        printf( "Lebre vence. Marmelada.\n" );
    }
    printf( "Total de batidas: %d\n", ticks );

    return 0;
}

void moveTartaruga( int *posicao )
{
    int i = 1 + rand() % 10;
    if      ( i >= 1 && i <= 5 ) *posicao += 3;   /* caminha rapido */
    else if ( i == 6 || i == 7 ) *posicao -= 6;   /* escorrega      */
    else                          *posicao += 1;   /* caminha devagar*/

    if ( *posicao < 1 ) *posicao = 1;
}

void moveLebre( int *posicao )
{
    int i = 1 + rand() % 10;
    if      ( i >= 1 && i <= 2 ) ;                 /* dorme: nao move */
    else if ( i >= 3 && i <= 4 ) *posicao += 9;    /* grande salto    */
    else if ( i == 5 )           *posicao -= 12;   /* escorregao      */
    else if ( i >= 6 && i <= 8 ) *posicao += 1;    /* pulinho         */
    else                         *posicao -= 2;    /* escorrega pouco */

    if ( *posicao < 1 ) *posicao = 1;
}

void exibeCorrida( int tartaruga, int lebre )
{
    int i;
    for ( i = 1; i <= PISTA; ++i ) {
        if ( i == tartaruga && i == lebre ) printf( "AI" );
        else if ( i == tartaruga ) printf( "T" );
        else if ( i == lebre ) printf( "L" );
        else printf( " " );
    }
    printf( "\n" );
}
\end{lstlisting}

\textbf{Conceito-chave do Cap.~7:} usamos \texttt{*posicao} dentro de \texttt{moveTartaruga}
e \texttt{moveLebre} para alterar a variável original. Essa é a \textbf{chamada por
referência simulada} -- exatamente o caso de uso central de ponteiros!
\end{respostabox}

\newpage

% ============================================================
\section{Exercício 7.18 -- Embaralhar + Distribuir em uma função}
% ============================================================

\begin{respostabox}
Combina-se a estrutura aninhada da \texttt{shuffle} com a lógica da \texttt{deal}
em uma única função que percorre cada carta uma vez, embaralhando enquanto
distribui. Essa abordagem é mais elegante e evita o adiamento indefinido.

A ideia central: ao percorrer o array sequencialmente após embaralhar (Fisher-Yates),
a posição \textit{é} a ordem natural de distribuição. As primeiras 5 são a mão 1,
as próximas 5 são a mão 2, e assim por diante.
\end{respostabox}

% ============================================================
\section{Exercício 7.19 -- O que \texttt{mystery1} faz?}
% ============================================================

\begin{respostabox}
\begin{lstlisting}
void mystery1( char *s1, const char *s2 )
{
    while ( *s1 != '\0' ) {  /* primeira fase: avancar ate o fim de s1 */
        s1++;
    }
    for ( ; *s1 = *s2; s1++, s2++ ) {
        ;  /* corpo vazio */
    }
}
\end{lstlisting}

\textbf{Análise:}
\begin{enumerate}[noitemsep]
  \item Primeira fase: o ponteiro \texttt{s1} avança até atingir o \texttt{'\textbackslash 0'}
  (terminador da string).
  \item Segunda fase: copia caractere por caractere de \texttt{s2} para o local
  apontado por \texttt{s1}, incluindo o \texttt{'\textbackslash 0'} de \texttt{s2}.
\end{enumerate}

\textbf{Conclusão:} \texttt{mystery1} é uma implementação manual de \textbf{\texttt{strcat}} --
\textit{concatena} \texttt{s2} ao final de \texttt{s1}.

\textbf{Exemplo:} se \texttt{string1 = "Ola"} e \texttt{string2 = "Mundo"}, ao final
\texttt{string1 = "OlaMundo"}.

\textbf{Saída final:} \texttt{OlaMundo}.

\textit{Nota:} a expressão \texttt{*s1 = *s2} na condição do \texttt{for} é uma
atribuição (não comparação!) que avalia para o caractere copiado. Quando o
\texttt{'\textbackslash 0'} é copiado, a condição se torna falsa e o loop termina.
\end{respostabox}

% ============================================================
\section{Exercício 7.20 -- O que \texttt{mystery2} faz?}
% ============================================================

\begin{respostabox}
\begin{lstlisting}
int mystery2( const char *s )
{
    int x;
    for ( x = 0; *s != '\0'; s++ ) {
        x++;
    }
    return x;
}
\end{lstlisting}

\textbf{Análise:} A função percorre cada caractere de \texttt{s} até encontrar
\texttt{'\textbackslash 0'}, incrementando \texttt{x} a cada caractere.

\textbf{Conclusão:} \texttt{mystery2} é uma implementação manual de \textbf{\texttt{strlen}} --
retorna o \textit{comprimento da string} (número de caracteres antes do
\texttt{'\textbackslash 0'}).

\textbf{Exemplo:} \texttt{mystery2("Hello")} retorna 5.
\end{respostabox}

\newpage

% ============================================================
\section{Exercício 7.21 -- Erros}
% ============================================================

\begin{respostabox}
\begin{enumerate}[label=\textbf{\alph*)}]

\item \textbf{Código:}
\begin{lstlisting}
int *number;
printf( "%d\n", *number );
\end{lstlisting}
\textbf{Erro:} ponteiro não inicializado, desreferenciado.
\textbf{Correção:} inicializar antes:
\begin{lstlisting}
int x = 42;
int *number = &x;
printf( "%d\n", *number );
\end{lstlisting}

\item \textbf{Código:}
\begin{lstlisting}
float *realPtr;
long *integerPtr;
integerPtr = realPtr;
\end{lstlisting}
\textbf{Erro:} atribuição entre ponteiros de tipos incompatíveis.
\textbf{Correção:} usar cast (mas verificar se faz sentido):
\begin{lstlisting}
integerPtr = ( long * )realPtr;
\end{lstlisting}

\item \textbf{Código:}
\begin{lstlisting}
int * x, y;
x = y;
\end{lstlisting}
\textbf{Erro:} \texttt{int * x, y;} declara \texttt{x} como \texttt{int *} e
\texttt{y} como \texttt{int} simples. Atribuir um \texttt{int} (\texttt{y}) a um
ponteiro (\texttt{x}) é inválido.
\textbf{Correção:}
\begin{lstlisting}
int *x, *y;  /* ambos sao ponteiros */
x = y;       /* OK */
\end{lstlisting}
Ou, se a intenção era atribuir endereço de \texttt{y} a \texttt{x}:
\begin{lstlisting}
int *x;
int y;
x = &y;
\end{lstlisting}

\item \textbf{Código:}
\begin{lstlisting}
char s[] = "este e um array de caracteres";
int count;
for ( ; *s != '\0'; s++ )
    printf( "%c ", *s );
\end{lstlisting}
\textbf{Erro:} \texttt{s} é nome de array -- não pode ser incrementado.
\textbf{Correção:} usar ponteiro auxiliar:
\begin{lstlisting}
char *p = s;
for ( ; *p != '\0'; p++ )
    printf( "%c ", *p );
\end{lstlisting}

\item \textbf{Código:}
\begin{lstlisting}
short *numPtr, result;
void *genericPtr = numPtr;
result = *genericPtr + 7;
\end{lstlisting}
\textbf{Erro:} desreferência de \texttt{void *} sem cast.
\textbf{Correção:}
\begin{lstlisting}
result = *( ( short * )genericPtr ) + 7;
\end{lstlisting}

\item \textbf{Código:}
\begin{lstlisting}
float x = 19.34;
float xPtr = &x;
printf( "%f\n", xPtr );
\end{lstlisting}
\textbf{Erros:} (1) \texttt{xPtr} declarado como \texttt{float} (deveria ser \texttt{float *});
(2) imprime ponteiro com \texttt{\%f}.
\textbf{Correção:}
\begin{lstlisting}
float *xPtr = &x;
printf( "%f\n", *xPtr );  /* ou usar %p para o endereco */
\end{lstlisting}

\item \textbf{Código:}
\begin{lstlisting}
char *s;
printf( "%s\n", s );
\end{lstlisting}
\textbf{Erro:} \texttt{s} é ponteiro não inicializado.
\textbf{Correção:}
\begin{lstlisting}
char *s = "Hello, World!";
printf( "%s\n", s );
\end{lstlisting}
\end{enumerate}
\end{respostabox}

\newpage

% ============================================================
\section{Exercício 7.22 -- Travessia de Labirinto}
% ============================================================

\begin{respostabox}
\textbf{Algoritmo da mão direita:} sempre seguir a parede da direita garante a saída
de qualquer labirinto simplesmente conexo. Implementação recursiva:

\begin{lstlisting}
/* Ex7.22: mazeTraverse.c -- esqueleto */
#include <stdio.h>
#define LADO 12

/* Direcoes: 0=Norte, 1=Leste, 2=Sul, 3=Oeste */
int dRow[] = { -1, 0, 1, 0 };
int dCol[] = {  0, 1, 0, -1 };

int mazeTraverse( char m[][ LADO ], int row, int col, int dir );
void exibeLab( char m[][ LADO ] );

int mazeTraverse( char m[][ LADO ], int row, int col, int dir )
{
    m[ row ][ col ] = 'X';
    exibeLab( m );

    /* Tenta primeiro a direita, depois a frente, depois esquerda */
    int tentativas[] = {
        ( dir + 1 ) % 4,   /* direita */
        dir,               /* frente  */
        ( dir + 3 ) % 4,   /* esquerda*/
        ( dir + 2 ) % 4    /* tras    */
    };

    int t, novaR, novaC, novaD;
    for ( t = 0; t < 4; ++t ) {
        novaD = tentativas[ t ];
        novaR = row + dRow[ novaD ];
        novaC = col + dCol[ novaD ];

        if ( novaR < 0 || novaR >= LADO ||
             novaC < 0 || novaC >= LADO ) {
            /* Saiu do labirinto = encontrou a saida! */
            return 1;
        }

        if ( m[ novaR ][ novaC ] == '.' ) {
            if ( mazeTraverse( m, novaR, novaC, novaD ) )
                return 1;
        }
    }

    /* Beco sem saida: desmarca e retrocede */
    m[ row ][ col ] = '.';
    return 0;
}
\end{lstlisting}
\end{respostabox}

% ============================================================
\section{Exercício 7.23 -- Gerar Labirintos}
% ============================================================

\begin{respostabox}
A geração aleatória de labirintos é um problema complexo. Algoritmos clássicos:
\begin{itemize}[noitemsep]
  \item \textbf{DFS recursivo:} começar com tudo parede, escolher célula inicial,
  explorar vizinhos não-visitados aleatoriamente, removendo paredes.
  \item \textbf{Algoritmo de Prim:} similar, mas escolhe sempre uma parede aleatória
  da fronteira.
  \item \textbf{Algoritmo de Kruskal:} usa \textit{union-find}.
\end{itemize}

A função \texttt{mazeGenerator} deve fornecer também as coordenadas de entrada e
saída do labirinto, geralmente nos extremos do tabuleiro.
\end{respostabox}

% ============================================================
\section{Exercício 7.24 -- Labirintos de qualquer tamanho}
% ============================================================

\begin{respostabox}
Para generalizar, basta passar o tamanho como parâmetro:
\begin{lstlisting}
int mazeTraverse( char **m, int linhas, int colunas,
                  int row, int col, int dir );
\end{lstlisting}
Em C, arrays bidimensionais de tamanho variável podem ser tratados como
\textbf{ponteiros para ponteiros} (\texttt{char **}) ou alocados dinamicamente
com \texttt{malloc} (Cap.~12).
\end{respostabox}

\newpage

% ============================================================
\section{Exercício 7.25 -- Arrays de Ponteiros para Funções}
% ============================================================

\begin{respostabox}
\begin{lstlisting}
/* Ex7.25: menu_grades.c
   Menu com array de ponteiros para funcoes */
#include <stdio.h>
#define ALUNOS 3
#define EXAMES 4

void printArray( const int grades[][ EXAMES ] );
void minimum   ( const int grades[][ EXAMES ] );
void maximum   ( const int grades[][ EXAMES ] );
void average   ( const int grades[][ EXAMES ] );

int main( void )
{
    /* Array de ponteiros para funcoes */
    void ( *processGrades[ 4 ] )( const int [][ EXAMES ] ) = {
        printArray, minimum, maximum, average
    };

    int grades[ ALUNOS ][ EXAMES ] = {
        { 77, 68, 86, 73 },
        { 96, 87, 89, 81 },
        { 70, 90, 86, 81 }
    };

    int escolha;

    do {
        printf( "\nDigite uma escolha:\n" );
        printf( " 0 Imprime o array de notas\n" );
        printf( " 1 Acha a menor nota\n" );
        printf( " 2 Acha a maior nota\n" );
        printf( " 3 Imprime a media de cada aluno\n" );
        printf( " 4 Encerra o programa\n" );
        printf( "Escolha: " );
        scanf( "%d", &escolha );

        if ( escolha >= 0 && escolha < 4 ) {
            ( *processGrades[ escolha ] )( grades );
        }
    } while ( escolha != 4 );

    return 0;
}

void printArray( const int grades[][ EXAMES ] )
{
    int aluno, exame;
    for ( aluno = 0; aluno < ALUNOS; ++aluno ) {
        for ( exame = 0; exame < EXAMES; ++exame ) {
            printf( "%4d ", grades[ aluno ][ exame ] );
        }
        printf( "\n" );
    }
}

void minimum( const int grades[][ EXAMES ] )
{
    int aluno, exame, baixa = 100;
    for ( aluno = 0; aluno < ALUNOS; ++aluno ) {
        for ( exame = 0; exame < EXAMES; ++exame ) {
            if ( grades[ aluno ][ exame ] < baixa )
                baixa = grades[ aluno ][ exame ];
        }
    }
    printf( "Menor nota: %d\n", baixa );
}

void maximum( const int grades[][ EXAMES ] )
{
    int aluno, exame, alta = 0;
    for ( aluno = 0; aluno < ALUNOS; ++aluno ) {
        for ( exame = 0; exame < EXAMES; ++exame ) {
            if ( grades[ aluno ][ exame ] > alta )
                alta = grades[ aluno ][ exame ];
        }
    }
    printf( "Maior nota: %d\n", alta );
}

void average( const int grades[][ EXAMES ] )
{
    int aluno, exame, soma;
    for ( aluno = 0; aluno < ALUNOS; ++aluno ) {
        soma = 0;
        for ( exame = 0; exame < EXAMES; ++exame ) {
            soma += grades[ aluno ][ exame ];
        }
        printf( "Aluno %d: media = %.2f\n",
                aluno + 1, ( double )soma / EXAMES );
    }
}
\end{lstlisting}

\textbf{Chave do exercício:} \texttt{processGrades[i]} contém um ponteiro para
uma das quatro funções. A escolha do usuário (0, 1, 2, 3) seleciona qual chamar --
substituindo um \texttt{switch} extenso por uma única linha:
\texttt{(*processGrades[escolha])(grades)}.
\end{respostabox}

\newpage

% ============================================================
\section{Exercício 7.26 -- O que \texttt{mystery3} faz?}
% ============================================================

\begin{respostabox}
\begin{lstlisting}
int mystery3( const char *s1, const char *s2 )
{
    for ( ; *s1 != '\0' && *s2 != '\0'; s1++, s2++ ) {
        if ( *s1 != *s2 ) {
            return 0;
        }
    }
    return 1;
}
\end{lstlisting}

\textbf{Análise:} A função compara caracteres correspondentes de \texttt{s1} e
\texttt{s2}, terminando ao primeiro \texttt{'\textbackslash 0'} encontrado em
qualquer das strings. Se algum par for diferente, retorna 0. Caso contrário,
retorna 1.

\textbf{Conclusão:} \texttt{mystery3} verifica se duas strings são \textbf{iguais}
(análoga à parte ``de igualdade'' de \texttt{strcmp}). \textbf{Atenção:} se uma
string é prefixo da outra (ex.: ``casa'' e ``casamento''), a função retorna 1
incorretamente, porque o loop termina assim que uma das strings acaba.

\textbf{Versão corrigida:}
\begin{lstlisting}
return *s1 == *s2;  /* verifica que ambos chegaram em \0 */
\end{lstlisting}
\end{respostabox}

% ============================================================
\section{Exercícios 7.27 a 7.29 -- Simpletron}
% ============================================================

\begin{respostabox}
\textbf{Visão geral:} estes três exercícios formam um projeto integrado: construir
um simulador completo de uma máquina virtual chamada \textbf{Simpletron}, com sua
própria linguagem de máquina (SML).

\textbf{Componentes do Simpletron:}
\begin{itemize}[noitemsep]
  \item \textbf{Memória:} 100 palavras de 4 dígitos com sinal.
  \item \textbf{Acumulador:} registrador especial para cálculos.
  \item \textbf{Contador de instrução:} aponta para próxima instrução.
  \item \textbf{Registrador de instrução:} contém a instrução atual.
\end{itemize}

\textbf{Conjunto de instruções SML:}
\begin{center}
\renewcommand{\arraystretch}{1.2}
\begin{tabular}{c l l}
\toprule
\textbf{Código} & \textbf{Nome} & \textbf{Função} \\ \midrule
10 & READ      & lê palavra do teclado \\
11 & WRITE     & imprime palavra \\
20 & LOAD      & carrega palavra no acumulador \\
21 & STORE     & armazena acumulador na memória \\
30 & ADD       & soma à palavra do acumulador \\
31 & SUBTRACT  & subtrai do acumulador \\
32 & DIVIDE    & divide o acumulador \\
33 & MULTIPLY  & multiplica o acumulador \\
40 & BRANCH    & desvio incondicional \\
41 & BRANCHNEG & desvio se acumulador negativo \\
42 & BRANCHZERO & desvio se acumulador zero \\
43 & HALT      & encerra execução \\
\bottomrule
\end{tabular}
\end{center}

\textbf{Exemplo SML (Exercício 7.27a -- soma de 10 inteiros positivos com sentinela):}
\begin{lstlisting}
00 +1015   /* READ no local 15      */
01 +2015   /* LOAD valor lido       */
02 +4112   /* BRANCHNEG para 12 (sentinela negativa = fim) */
03 +2014   /* LOAD soma              */
04 +3015   /* ADD valor              */
05 +2114   /* STORE soma             */
06 +4000   /* BRANCH para 00 (proximo) */
12 +1114   /* WRITE soma             */
13 +4300   /* HALT                   */
14 +0000   /* soma                   */
15 +0000   /* valor lido             */
\end{lstlisting}

\textbf{Estrutura do simulador (Ex.\ 7.28):}
\begin{lstlisting}
/* Ciclo de execucao da instrucao */
while ( !halt ) {
    instructionRegister = memory[ instructionCounter ];
    operationCode = instructionRegister / 100;
    operand       = instructionRegister % 100;

    switch ( operationCode ) {
        case READ:
            scanf( "%d", &memory[ operand ] );
            break;
        case LOAD:
            accumulator = memory[ operand ];
            break;
        case ADD:
            accumulator += memory[ operand ];
            break;
        case BRANCH:
            instructionCounter = operand;
            continue;  /* nao incrementa o IC */
        case HALT:
            halt = 1;
            break;
        /* ... demais casos ... */
    }
    ++instructionCounter;
}
\end{lstlisting}

\textbf{Reflexão:} Construir um simulador de máquina é uma das experiências mais
formativas em ciência da computação. Conecta hardware (ciclos de busca/execução),
arquitetura (organização de memória, registradores) e linguagens (interpretação
de códigos de operação). Os Exercícios 7.27--7.29 são extensos -- o livro dedica
páginas a eles -- mas o esforço vale muito a pena.
\end{respostabox}

\newpage

% ============================================================
\section{Exercício 7.30 -- Cálculos Geométricos com Ponteiros para Função}
% ============================================================

\begin{respostabox}
\begin{lstlisting}
/* Ex7.30: geometria_funcoes.c
   Menu com array de ponteiros para funcao */
#include <stdio.h>
#define PI 3.14159

void circunferencia( double raio );
void areaCirculo   ( double raio );
void volumeEsfera  ( double raio );

int main( void )
{
    /* Array de ponteiros para funcao */
    void ( *calculos[ 3 ] )( double ) = {
        circunferencia, areaCirculo, volumeEsfera
    };

    int escolha;
    double raio;

    do {
        printf( "\n0 - Circunferencia\n" );
        printf( "1 - Area do circulo\n" );
        printf( "2 - Volume da esfera\n" );
        printf( "3 - Encerrar\n" );
        printf( "Escolha: " );
        scanf( "%d", &escolha );

        if ( escolha >= 0 && escolha <= 2 ) {
            printf( "Digite o raio: " );
            scanf( "%lf", &raio );
            ( *calculos[ escolha ] )( raio );
        }
    } while ( escolha != 3 );

    return 0;
}

void circunferencia( double raio )
{
    printf( "Circunferencia (r=%.2f) = %.4f\n", raio, 2 * PI * raio );
}

void areaCirculo( double raio )
{
    printf( "Area (r=%.2f) = %.4f\n", raio, PI * raio * raio );
}

void volumeEsfera( double raio )
{
    printf( "Volume da esfera (r=%.2f) = %.4f\n",
            raio, ( 4.0 / 3.0 ) * PI * raio * raio * raio );
}
\end{lstlisting}
\end{respostabox}

% ============================================================
\section{Exercício 7.31 -- Cálculos com Dois Operandos}
% ============================================================

\begin{respostabox}
\begin{lstlisting}
/* Ex7.31: calculadora.c */
#include <stdio.h>

void somar      ( double a, double b );
void subtrair   ( double a, double b );
void multiplicar( double a, double b );
void dividir    ( double a, double b );

int main( void )
{
    void ( *operacoes[ 4 ] )( double, double ) = {
        somar, subtrair, multiplicar, dividir
    };

    int escolha;
    double a, b;

    do {
        printf( "\n0-Somar 1-Subtrair 2-Multiplicar 3-Dividir 4-Sair\n" );
        printf( "Escolha: " );
        scanf( "%d", &escolha );

        if ( escolha >= 0 && escolha <= 3 ) {
            printf( "Digite dois numeros: " );
            scanf( "%lf%lf", &a, &b );
            ( *operacoes[ escolha ] )( a, b );
        }
    } while ( escolha != 4 );

    return 0;
}

void somar( double a, double b )
{
    printf( "%.2f + %.2f = %.2f\n", a, b, a + b );
}
void subtrair( double a, double b )
{
    printf( "%.2f - %.2f = %.2f\n", a, b, a - b );
}
void multiplicar( double a, double b )
{
    printf( "%.2f * %.2f = %.2f\n", a, b, a * b );
}
void dividir( double a, double b )
{
    if ( b == 0.0 ) {
        printf( "Erro: divisao por zero!\n" );
    }
    else {
        printf( "%.2f / %.2f = %.2f\n", a, b, a / b );
    }
}
\end{lstlisting}
\end{respostabox}

\end{document}
